\documentclass{amsart}
\usepackage{mathrsfs}
\usepackage{quiver}

\usepackage{amsmath,amssymb,amsfonts,amscd}
\usepackage[all,cmtip]{xy}
\usepackage{enumerate}
\usepackage{ amssymb, latexsym, amsmath}
%\setcounter{MaxMatrixCols}{30}%
\usepackage{fullpage}
\usepackage{url}
\usepackage{hyperref}
\usepackage{ marvosym }




% \numberwithin{equation}{section}

%\usepackage{cjw-latex}

%\providecommand{\U}[1]{\protect\rule{.1in}{.1in}}

\theoremstyle{plain}
% \newtheorem{theorem}{Theorem}[section]
\newtheorem{theorem}{Theorem}
\newtheorem*{thm}{Theorem}
\newtheorem{proposition}[theorem]{Proposition}
\newtheorem{lemma}[theorem]{Lemma}
\newtheorem{conjecture}[theorem]{Conjecture}
\newtheorem{corollary}[theorem]{Corollary}
\newtheorem{algorithm}[theorem]{Algorithm}
\newtheorem{axiom}[theorem]{Axiom}
\newtheorem{criterion}[theorem]{Criterion}
%\newtheorem{conjecture}[theorem]{Conjecture}
\newtheorem{wildconjecture}[theorem]{Wild Conjecture}


\theoremstyle{definition}
\newtheorem{definition}[theorem]{Definition}
\newtheorem*{dfn}{Definition}
\newtheorem{condition}[theorem]{Condition}
\newtheorem{example}[theorem]{Example}
\newtheorem{exercise}[theorem]{Exercise}
\newtheorem{notation}[theorem]{Notation}
\newtheorem{question}[theorem]{Question}
\newtheorem{problem}[theorem]{Problem}
\newtheorem{solution}[theorem]{Solution}




\theoremstyle{remark}
\newtheorem{remark}[theorem]{Remark}
\newtheorem{remarks}[theorem]{Remarks}
\newtheorem{summary}[theorem]{Summary}
\newtheorem{observation}[theorem]{Observation}
\newtheorem{conclusion}[theorem]{Conclusion}
\newtheorem{acknowledgement}[theorem]{Acknowledgement}
\newtheorem{case}[theorem]{Case}
\newtheorem{claim}[theorem]{Claim}



\makeatletter
\newcommand{\rmnum}[1]{\romannumeral #1}
\newcommand{\Rmnum}[1]{\expandafter\@slowromancap\romannumeral #1@}
\makeatother
\newcommand{\Aut}{\operatorname{Aut}}
\newcommand{\Ext}{\operatorname{Ext}}
\newcommand{\Tor}{\operatorname{Tor}}
\newcommand{\Z}{\mathbb{Z}}
\newcommand{\osum}{\oplus}
\newcommand{\aff}{\operatorname{Aff}}

\newcommand{\Id}{\operatorname{Id}}
\newcommand{\concat}{(\Id \to \Omega \Sigma)}
\newcommand{\eval}{(\Sigma\Omega \to \Id)}
\newcommand{\Q}{\mathbb{Q}}
\newcommand{\F}{\mathbb{F}}
\newcommand{\N}{\mathbb{N}}
\newcommand{\g}{\mathfrak{g}}
\newcommand{\n}{\mathfrak{n}}
\newcommand{\h}{\mathfrak{h}}
\newcommand{\p}{\mathfrak{p}}
\newcommand{\q}{\mathfrak{q}}
\newcommand{\m}{\mathfrak{m}}
\newcommand{\e}{\mathfrak{e}}
\newcommand{\f}{\mathfrak{f}}
\newcommand{\V}{\mathscr{V}}
\newcommand{\D}{\mathscr{D}}
\renewcommand{\S}{\mathscr{S}}
\newcommand{\Sch}{\mathscr{C}}
\newcommand{\Y}{\mathscr{Y}}
\renewcommand{\u}{\mathfrak{u}}
\newcommand{\set}[2]{ \left\{ {#1} \, \left| \, {#2} \right\}\right.}
\newcommand{\A}{\mathbb{A}}
\newcommand{\X}{\mathcal{X}}
\renewcommand{\Im}{\operatorname{Im}}
\newcommand{\T}{\mathscr{T}}
\newcommand{\HH}{\mathscr{H}}
\newcommand{\PGL}{\operatorname{PGL}}
\newcommand{\righthookarrow}{\hookrightarrow}
\newcommand{\sign}{\operatorname{sign}}
\renewcommand{\_}[2]{\underbrace{#1}_{#2}}
\renewcommand{\^}[2]{\overbrace{#1}_{#2}}
\newcommand{\curverightarrow}{\curvearrowright}

\newcommand{\z}{\mathfrak{z}}
\renewcommand{\sl}{\mathfrak{sl}}
\newcommand{\R}{\mathbb{R}}
\newcommand{\RP}{\mathbb{RP}}
\renewcommand{\P}{\mathbb{P}}
\newcommand{\C}{\mathbb{C}}
\renewcommand{\H}{\mathcal{H}}
\newcommand{\Ind}{\operatorname{Ind}}
\newcommand{\Tr}{\operatorname{Tr}}
\newcommand{\Orb}{\mathcal{O}}
\renewcommand{\k}{\mbox{\Fontauri k}}
\newcommand{\Ad}{\operatorname{Ad}}
\newcommand{\ad}{\operatorname{ad}}
\newcommand{\Hom}{\operatorname{Hom}}
\newcommand{\Ker}{\operatorname{Ker}}
\newcommand{\End}{\operatorname{End}}
\newcommand{\supp}{\operatorname{supp}}
\newcommand{\Stab}{\mbox{Stab}}
\newcommand{\Lie}{\operatorname{Lie}}
\newcommand{\GL}{\operatorname{GL}}
\newcommand{\SL}{\operatorname{SL}}
\newcommand{\Sp}{\operatorname{Sp}}
\newcommand{\Mp}{\operatorname{Mp}}
\renewcommand{\a}{\mathfrak{a}}
\newcommand{\Wh}{\operatorname{Wh}}
\newcommand{\Span}{\operatorname{Span}}
\newcommand{\WF}{\operatorname{WF}}
\newcommand{\AC}{\operatorname{AC}}
\newcommand{\Lin}{\operatorname{Lin}}
\newcommand{\Diff}{\operatorname{Diff}}
\newcommand{\Gen}{\operatorname{Gen\, Hom}}
\newcommand{\Her}{\operatorname{Her}}
\newcommand{\pivot}{&}
\newcommand{\sgn}{\operatorname{sgn}}
%\newcommand{C^{\infty}}{DS}
\newcommand{\vsp}{{\vspace{0.2in}}}
\newcommand{\cU}{\mathcal{U}}
\newcommand{\cO}{\mathcal{O}}
\newcommand{\sF}{\mathscr{F}}
\newcommand{\cB}{\mathcal{B}}

\title{Notes}

\author{George}
\begin{document}
\maketitle

\begin{theorem}
    If $X$ is a topological space and $\cU = \{U_\lambda \}_{\lambda \in \Lambda}$ then there is a bijection
    \begin{equation}
        \left\{
            \begin{gathered}
            \textnormal{Locally free sheaves on } X\\
            \textnormal{trivialised over } \cU
            \end{gathered}
        \right\}/\cong \quad\longleftrightarrow\quad \check H {(\cU, \GL_n(\cO_X))}
    \end{equation}

\end{theorem}

\begin{proof}
    Let us call the $\rightarrow$ direction $\alpha$ and the $\leftarrow$ direction $\beta$. We first construct $\alpha$ Let $\sF$ be a locally free sheaf on $X$ of rank $n$. Then for each $\lambda \in \Lambda$ there exists an ``trivialisation'' isomorphism of sheaves $\phi_\lambda: \cO_U^{\oplus n} \to \sF|_{U_\lambda}$. Let $U_{\lambda \mu} = U_\lambda \cap U_\mu$. Since $\phi_\lambda$ is a morphism of sheaves, there is a restriction $\phi_{\lambda}|_{U_{\lambda,\mu}}: \cO_U^{\oplus n}(U_{\lambda\mu}) \to \sF|_{U_{\lambda}}(U_{\lambda\mu})$. We can also restrict $\phi_\mu$ and the source and target can be rewritten as in the following diagram, 
    % https://q.uiver.app/#q=WzAsMixbMCwwLCJcXGNPX1hee1xcb3BsdXMgbn0oVV97XFxsYW1iZGEsXFxtdX0pIl0sWzIsMCwiXFxzRihVX3tcXG11LFxcbGFtYmRhfSkiXSxbMCwxLCJcXHBoaV97XFxtdX18X3tVX3tcXGxhbWJkYSxcXG11fX0iLDAseyJjdXJ2ZSI6Mn1dLFswLDEsIlxccGhpX3tcXGxhbWJkYX18X3tVX3tcXGxhbWJkYSxcXG11fX0iLDAseyJjdXJ2ZSI6LTJ9XV0=
\[\begin{tikzcd}
	{\cO_X^{\oplus n}(U_{\lambda,\mu})} && {\sF(U_{\mu,\lambda})}
	\arrow["{\phi_{\mu}|_{U_{\lambda,\mu}}}", curve={height=12pt}, from=1-1, to=1-3]
	\arrow["{\phi_{\lambda}|_{U_{\lambda,\mu}}}", curve={height=-12pt}, from=1-1, to=1-3]
\end{tikzcd}\]
\\
Now we can define $T_{\lambda \mu} = \phi_{\lambda}|_{U_{\lambda,\mu}}^{-1} \circ \phi_{\mu}|_{U_{\lambda\mu}} \in \Aut(\cO_X^{\oplus n}(U_{\lambda\mu})) = \GL_n(\cO_X)(U_{\lambda\mu})$. We now verify that $(T_{\lambda\mu} )_{\lambda,\mu \in \Lambda}$ satisfy the cocycle condition. Restricting to $U_{\lambda \mu \nu}$, we have
\begin{equation}
    T_{ \nu \mu}T_{ \mu \lambda} = (\phi_\nu^{-1} \phi_\mu) (\phi_{\mu}^{-1} \phi_{\lambda}) = \phi_\nu^{-1}\phi_{\lambda} = T_{\nu \lambda}.
\end{equation}
Therefore $(T_{\lambda \mu})_{\lambda,\mu} \in \check Z(\cU, \GL_n(\cO_X))$.
\\
Now suppose that $\sF'$ is another locally free sheaf trivialised over $\cU$, with $\phi_\lambda'$ and $T_{\lambda \mu}$ defined analogously, and suppose $\psi$ is an isomorphism of sheaves $\sF \to \sF'$. Then for each $\lambda \in \Lambda$ we have an isomorphism of sheaves $\psi_\lambda: \sF|_{U_\lambda} \to \sF'|_{U_\lambda}$. Let $h_\lambda = (\phi'_\lambda)^{-1} \circ \psi_\lambda \circ \phi_\lambda$. Clearly $h_\lambda \in \Gamma(U_\lambda,\GL_n(\cO_X))$, so $(h_\lambda)_{\lambda \in \Lambda} \in \cB^0(\cU, \GL_n(\cO_X))$. Restricting to $U_{\lambda \mu}$, since $\psi$ is a morphism of sheaves, $\psi_\lambda = \psi_\mu$, so
\begin{align}
    h_\lambda T_{\lambda \mu} &= ((\phi_\lambda')^{-1} \psi_\lambda \phi_\lambda) (\phi_\lambda^{-1} \phi|_\mu) \\
    &= (\phi_\lambda')^{-1} \psi_\lambda \phi_\mu = (\phi|_\lambda')^{-1} \psi|_\mu \phi_\mu \\
    &= (\phi_\lambda')^{-1} (\phi_\mu' (\phi|_\mu')^{-1}) \psi|_\mu \phi_\mu \\
    &= T_{\lambda \mu}' h_{\mu}.
\end{align}
\end{proof}
Therefore $(T_{\lambda\mu})$ and $(T_{\lambda\mu}')$ are in the same class of $\check H(\cU, \GL_n(\cO_X))$. Thus we have a well defined map $\alpha$. \\

Suppose that there exists $(h_\lambda)_{\lambda \in \Lambda} \in \cB^0(\cU, \GL_n(\cO_X))$ such that for each $\lambda, \mu \in \Lambda$, $h_\lambda T_{\lambda \mu} = T_{\lambda \mu}' h_{\mu}$ on $U_{\lambda \mu}$. Then let $\psi_\lambda = \phi'_\lambda h_\lambda \phi_\lambda^{-1}$. Then on each overlap $U_{\lambda\mu}$, we have 
\begin{align}
    \psi_\lambda &= \phi'_\lambda h_\lambda \phi_\lambda^{-1} \\
    &= \phi'_\lambda h_\lambda \phi_\lambda^{-1} (\phi_\mu \phi_\mu^{-1}) \\
    &= \phi'_\lambda h_\lambda T_{\lambda\mu} \phi_\mu^{-1} \\
    &= \phi'_\lambda T_{\lambda\mu}' h_\mu \phi_\mu^{-1} \\
    &= \psi_\mu.
\end{align}
Therefore the $\psi_\lambda$ agree on intersections, so that induce a morphism sheaves, which is in fact an isomorphism because the same argument holds for $(h_\lambda^{-1})$. Thus $\alpha$ is injective. \\

Now for the reverse direction, suppose we have a cocyle $(T_{\lambda\mu})_{\lambda,\mu \in \Lambda} \in \check Z(\cU, \GL_n(\cO_X))$. We first construct a sheaf $\sF$ such that $\sF(U_\lambda) = \cO_X^{\oplus n}(U_\lambda)$ for each $\lambda \in \Lambda$. For any open set $V$, we define
\begin{equation}
    \sF(V) = \left\{ (s_\lambda)_{\lambda \in \Lambda} \in \prod_{\lambda\in \Lambda} \cO_X^{\oplus n}(V \cap U_\lambda) \mid s_\lambda = T_{\lambda \mu} s_\mu \text{ on } V\cap U_\lambda \cap U_\mu\right\}.
\end{equation}
The restriction maps $\sF (V) \to \sF(W)$ for $V \subset W$ are induced by restricting $\cO_X(V \cap U_\lambda) \to \cO_X(W \cap U_\lambda)$ for each $\lambda$, and are obviously well defined, so $\sF$ is a presheaf. The separation and gluing axioms are satisfied fro $\sF$ as a direct consequence of the fact that they are satisfied for $\cO_X$, and a section of $\sF(V)$ is just a collection of sections of $\cO_X(V \cap U_\lambda)$. Thus $\sF$ is a sheaf. \\

For any open $V$, to give a section of $\sF(V)$ it suffices to give a section $s_\lambda$ for some $\lambda \in \Lambda$ such that $V \subseteq U_\lambda$, since we can then define $s_\lambda = T_{\lambda \mu} s_\mu$ for each $\mu \in \Lambda$. It follows that the inclusion into the $\lambda$ component gives an isomorphism $\phi_\lambda: \cO_X^{\oplus n}(U_\lambda) \to \sF(U_\lambda)$ for each $\lambda \in \Lambda$. Naturally, $\phi_\lambda^{-1}$ is the projection onto the $\lambda$ component. Therefore the transition functions of $\sF$ are $\phi_{\mu}^{-1} \phi_{\lambda} = T_{\lambda \mu} $. Thus we have a well defined map $\beta$, and $\alpha \circ \beta = \Id$, which means that $\alpha$ is surjective. We already proved that $\alpha$ is injective, so this completes the proof of the bijection.


\end{document}
